%==============================================================================
% Capítulo 4: Marco Metodológico
%==============================================================================
\section{Marco metodológico}\label{sec:marco-metodologico}

\subsection{Enfoque y tipo de investigación}\label{subsec:enfoque-investigacion}

La investigación se enmarca dentro del enfoque cuantitativo aplicado, dado que propone el desarrollo de una solución informática orientada a mejorar el funcionamiento de un proceso organizacional concreto y mide su impacto mediante indicadores numéricos. \textcite{HernandezSampieri2018_metodologia} ubican este tipo de estudios dentro de la investigación tecnológica de carácter aplicado, donde el conocimiento producido se orienta a la transformación práctica de la realidad antes que a la generación de teoría general. El trabajo posee adicionalmente carácter experimental, en la medida en que diseña, implementa y evalúa un artefacto tecnológico bajo condiciones controladas.

\subsection{Diseño experimental}\label{subsec:diseno-experimental}

El diseño experimental adoptado es de tipo cuasiexperimental con muestras pareadas y mediciones repetidas. La unidad de observación es el incidente individual, y cada uno se procesa secuencialmente bajo dos condiciones: la condición manual y la condición automatizada. Esta estrategia ---procesar el mismo conjunto de incidentes por ambas vías--- reduce sustancialmente la variabilidad asociada a las características intrínsecas de cada caso, lo que aumenta la potencia estadística del contraste entre flujos.

Para acotar el efecto de aprendizaje del operador humano sobre el corpus, el diseño adoptó una ejecución secuencial pareada: los mismos 200 casos se procesaron primero por el flujo manual (junio a agosto de 2026) y luego por el flujo automatizado (septiembre a octubre de 2026; \autoref{fig:cronograma}). El procesamiento manual se distribuyó entre dos operadores experimentados, rotados entre las dos mitades del corpus (100 casos por mitad). Con dos operadores y dos mitades, no obstante, la mitad del corpus y el operador quedan colineales: la rotación de operadores queda confundida con el efecto del operador y no permite estimar por separado el efecto de aprendizaje, de modo que no se presenta aquí como una mitigación efectiva de la amenaza interna sino como una limitación reconocida del diseño (véase la \autoref{subsec:amenazas-validez}). El diseño conserva, sí, el control sobre la variabilidad asociada a las características intrínsecas de cada caso, al medir el mismo incidente bajo ambas condiciones.

La investigación se desenvolvió bajo un protocolo de medición en vivo. Los dos operadores trabajaron en pares y se cronometraron mutuamente, de modo que sabían que la campaña registraba tiempos; no se anunciaron el inicio de cada cronometraje ni el propósito comparativo del estudio (el contraste con un flujo automatizado), informado al concluir la campaña en una sesión de esclarecimiento (\emph{debriefing}). Así se atenuó, sin eliminarlo, el efecto Hawthorne ---la alteración del comportamiento de los sujetos cuando saben que están siendo observados--- \parencite{HernandezSampieri2018_metodologia}, que subsiste como amenaza residual (\autoref{sec:discusion}). En el flujo manual, el compañero de par que no registraba el incidente cronometró en vivo su tiempo total, consignado en la planilla de cronometraje en segundos, con redondeo hacia arriba. De los 200 valores, 199 son múltiplos de 5~s y 180 de 10~s, resolución efectiva del instrumento que vincula los descuentos de 5~s y 10~s del análisis de sensibilidad con el procedimiento. El redondeo al alza constituye un sesgo de medición sistemático que infla el tiempo manual y su efecto se acota mediante el análisis de sensibilidad reportado en la \autoref{sec:resultados}.

\subsection{Operacionalización de variables}\label{subsec:operacionalizacion-variables}

La Tabla~\ref{tab:variables} presenta la operacionalización de las variables del estudio, distinguiendo entre variables cuantitativas continuas, cuantitativas proporcionales y cualitativas nominales. Cada variable se acompaña de su unidad de medida y de la definición operacional aplicada durante la captura, garantizando la trazabilidad entre la definición conceptual y el procedimiento de medición efectivo.

\begin{table}[htbp]
\centering
\caption{Operacionalización de variables del estudio}
\label{tab:variables}
\small
\begin{tabular}{@{}>{\raggedright\arraybackslash}p{2.5cm} >{\raggedright\arraybackslash}p{3.2cm} >{\raggedright\arraybackslash}p{2.3cm} >{\raggedright\arraybackslash}p{5.4cm}@{}}
\toprule
\textbf{Variable} & \textbf{Tipo} & \textbf{Unidad} & \textbf{Definición operacional} \\
\midrule
Tiempo de registro & Cuantitativa continua & Segundos (s) & Lapso entre la recepción del incidente en el canal de entrada y la persistencia confirmada del ticket en la base de datos. \\
Exactitud de clasificación & Cuantitativa proporcional & Porcentaje [0; 100] & Cociente entre incidentes correctamente clasificados y total de incidentes evaluados. \\
F1 macro promediado & Cuantitativa proporcional & Adimensional [0; 1] & Promedio aritmético del F1 de cada sector; el de umbral excluye los de soporte unitario. \\
Intervención humana & Cuantitativa proporcional & Porcentaje [0; 100] & Se reporta de dos formas rotuladas: proporción de casos que requieren acción manual de un operador humano y proporción del tiempo total del proceso dedicado a esa acción. \\
Canal de origen & Cualitativa nominal & Categoría & Vía de entrada: correo electrónico, formulario web o llamada telefónica. \\
Categoría asignada & Cualitativa nominal & Categoría & Sector responsable, uno de los cinco sectores canónicos: Seguridad Informática, Soporte Técnico Hardware, Soporte Técnico Software, Bases de Datos o Sistemas. \\
\bottomrule
\end{tabular}
\par\smallskip
{\footnotesize\textit{Fuente:} elaboración propia.\par}
\end{table}

\subsection{Población, muestra y corpus}\label{subsec:poblacion-muestra-corpus}

La población de referencia está constituida por la totalidad de los incidentes informáticos reportados a la mesa de ayuda de la organización analizada. A partir de un universo aproximado de 2.646 incidentes ---proyectado como 42 incidentes diarios sobre los 63 días hábiles del período relevado (\(42 \times 63 \approx 2.646\))---, se construyó, por muestreo estratificado con asignación proporcional a la distribución por sector observada en el período relevado, un corpus de 200 casos pseudonimizados distribuidos entre los cinco sectores canónicos: Soporte Técnico Hardware 85, Soporte Técnico Software 67, Seguridad Informática 31, Sistemas 16 y Bases de Datos 1. La distribución del universo por sector no está disponible, por lo que esa proporcionalidad no es verificable; Bases de Datos quedó con un único caso, no estimable por clase, y en replicaciones convendría sobremuestrear las clases minoritarias. Los casos provienen de tres canales de origen: llamada telefónica 81, correo electrónico 66 y formulario web 53. El flujo manual se cronometró entre junio y agosto de 2026; sobre los mismos 200 casos, el flujo automatizado se ejecutó entre septiembre y octubre de 2026.

El tamaño muestral se determinó mediante la fórmula clásica de muestreo aleatorio simple con corrección por finitud, apropiada porque el universo relevado es finito y relativamente acotado:
\[
n = \frac{N\, z^{2}\, p\,(1-p)}{e^{2}\,(N-1) + z^{2}\, p\,(1-p)},
\]
equivalente a \(n = n_0 / [1 + (n_0 - 1)/N]\) con \(n_0 = z^{2}\, p\,(1-p)/e^{2}\). En ella, \(N \approx 2.646\) es el tamaño del universo trimestral, \(z = 1{,}96\) corresponde a un nivel de confianza del 95~\%, \(p = 0{,}5\) es la proporción esperada conservadora (la que maximiza la varianza \(p(1-p)\) y, con ella, el tamaño requerido) y \(e\) es el margen de error absoluto. El tamaño del corpus se fijó, no obstante, por disponibilidad operativa: los 200 casos son los incidentes efectivamente cronometrados en la campaña manual, no un tamaño derivado a priori de la fórmula. Con esos insumos y ese tamaño fijo, el margen de error implicado bajo corrección por finitud es del 6,66~\%, precisión que el corpus construido admite; recíprocamente, con un margen del 6,66~\% la fórmula reproduce el tamaño de 200 casos (\(n \approx 200{,}0\)). El margen resultante es razonable para un estudio piloto orientado a evaluar viabilidad técnica antes de un despliegue definitivo.

La construcción del corpus aplicó criterios de exclusión explícitos. En primer lugar, se descartaron los incidentes duplicados, entendiendo por tales los reportados más de una vez sobre el mismo hecho, ya sea por repetición en un mismo canal o por ingreso concurrente por canales distintos. En segundo lugar, se descartaron las descripciones vacías o insuficientes, esto es, los registros sin contenido textual clasificable o cuya extensión no permitía inferir un sector responsable. En cambio, los casos con pertenencia simultánea a más de un sector (multicategoría) y las descripciones muy breves pero clasificables se conservaron deliberadamente, por ser representativos del entorno operativo real, asignándoseles una etiqueta primaria como sector responsable. Este criterio es coherente con el análisis de errores de la \autoref{sec:resultados}, donde las descripciones extremadamente cortas y la mezcla de categorías figuran entre los patrones de error dominantes: excluirlos habría sesgado el corpus hacia casos artificialmente limpios y sobreestimado el desempeño del clasificador.

Cada incidente del corpus fue revisado y etiquetado de forma independiente por dos analistas con experiencia en la mesa de ayuda, ajenos al equipo de autores del trabajo y ajenos también al diseño del clasificador evaluado y a la fijación de sus umbrales. La independencia que respalda la validez de constructo es de carácter funcional: los analistas no participaron en la construcción del sistema, no fueron instruidos según su comportamiento y no tuvieron acceso a las predicciones del clasificador al momento de etiquetar. Los dos analistas son, además, las mismas dos personas que operaron el brazo manual descrito en la \autoref{subsec:protocolo-pruebas}. Esta coincidencia se declara explícitamente: no altera el cómputo del acuerdo inter-evaluador, que mide la concordancia entre ambos etiquetadores y no su relación con las condiciones del experimento, pero implica que la etiqueta de referencia y el brazo de comparación comparten sujetos, lo que acota la independencia del \emph{ground truth} respecto de los operadores cronometrados y se retoma entre las amenazas a la validez de constructo (\autoref{subsec:amenazas-validez}). Las discrepancias entre analistas fueron resueltas por consenso bajo la supervisión del director del trabajo. El acuerdo entre evaluadores se calculó mediante el coeficiente Kappa de Cohen \parencite{Cohen1960_kappa}; el valor obtenido (\(\kappa = 0{,}87\)) se interpreta como casi perfecto según los rangos habitualmente aceptados, que sitúan el acuerdo casi perfecto por encima de \(\kappa > 0{,}80\) \parencite{Landis1977_kappa}. Este valor respalda la confiabilidad de la categoría de referencia utilizada como verdad fundamental (\emph{ground truth}) para evaluar el clasificador automático. No se reporta el intervalo de confianza del kappa: su cálculo exige la tabla de acuerdo observado/esperado o las etiquetas por caso de ambos analistas, registros que no se conservaron en las fuentes disponibles; el valor puntual se informa, por lo tanto, sin incertidumbre cuantificada, y el intervalo queda pendiente de esa tabla primaria.

\subsection{Protocolo de pruebas}\label{subsec:protocolo-pruebas}

El protocolo de pruebas se ejecutó en dos fases sucesivas. La primera fase consistió en una validación funcional del sistema, en la cual se verificó la correcta integración de los componentes mediante: (a)~pruebas unitarias del módulo Python, cuya cobertura se mide con \texttt{pytest-cov} sobre los paquetes de rutas, servicios y repositorios mediante el comando \texttt{pytest --cov=app.routes --cov=app.services --cov=app.repositories --cov-report=term-missing}, (b)~pruebas de integración de la API y de la base de datos, junto con pruebas estructurales que validan el archivo \texttt{n8n/workflow.json}, y (c)~pruebas de aceptación simuladas con incidentes sintéticos representativos de cada categoría.

La segunda fase consistió en la validación experimental propiamente dicha, en la cual los 200 incidentes del corpus se procesaron bajo ambas condiciones operativas y se registraron las variables descritas en la \autoref{subsec:operacionalizacion-variables}. El procesamiento manual fue realizado por dos operadores experimentados con más de tres años de antiguedad en la mesa de ayuda, ambos capacitados previamente en el uso del sistema de gestión de tickets utilizado como referencia. El procesamiento automatizado se ejecutó sobre la misma infraestructura, sin intervención humana en la captura y con el gate de revisión aplicable, durante la ventana de septiembre a octubre de 2026 y con monitoreo continuo de los tiempos de respuesta.

\subsection{Métricas e instrumentos}\label{subsec:metricas-instrumentos}

La evaluación del sistema se sustenta en cuatro métricas principales. La primera es el tiempo medio de registro, calculado como la media aritmética y la mediana de los tiempos individuales bajo cada condición, junto con su desvío estándar muestral, mínimo, máximo y reducción relativa de la media. La segunda es la exactitud global de clasificación, calculada como la proporción de incidentes correctamente derivados al sector responsable. La tercera es el valor F1 macro promediado, media aritmética de los valores F1 de cada clase estimable, que neutraliza el desbalanceo de clases del corpus \parencite{Sokolova2009_performance_measures}. La cuarta es la proporción de intervención humana, prevista en dos dimensiones rotuladas: la proporción de casos que requieren intervención y la proporción del tiempo dedicado a tareas manuales sobre el tiempo total del proceso; en la versión evaluada no se midió, por ausencia de registro primario.

Como instrumentos de captura se utilizaron: (a)~registros automáticos generados por el propio sistema con marcas de tiempo precisas al milisegundo para el flujo automatizado, (b)~una planilla de cronometraje observacional, completada en vivo por el compañero de par que no estaba registrando el incidente, que consigna el tiempo total de registro por incidente en segundos, con redondeo hacia arriba, para el flujo manual, y (c)~matrices de confusión generadas mediante la biblioteca scikit-learn \parencite{Pedregosa2011_scikit_learn}. En el flujo automatizado, el tiempo de registro del canal telefónico comprende tanto la duración de la llamada grabada como la latencia de procesamiento posterior al cuelgue, mientras que los canales de correo electrónico y formulario web se miden como latencia de extremo a extremo de su procesamiento. Dado que la grabación telefónica está acotada a un máximo de 45~s (véase la \autoref{sec:implementacion}), la duración de la llamada registrada ---y por lo tanto el tiempo automatizado telefónico--- constituye una cota inferior de la duración real de la llamada.

\subsection{Análisis estadístico}\label{subsec:analisis-estadistico}

Para la comparación de los tiempos de registro entre ambas condiciones se aplicó la prueba de Wilcoxon de rangos con signo \parencite{Wilcoxon1945_individual_comparisons}, prueba no paramétrica para muestras pareadas, dado que la prueba previa de normalidad de Shapiro-Wilk rechazó la hipótesis de normalidad en los tres conjuntos analizados, todos con \(n = 200\): tiempos manuales (\(W = 0{,}803\), \(p = 3{,}71 \times 10^{-15}\)), tiempos automatizados (\(W = 0{,}902\), \(p = 3{,}47 \times 10^{-10}\)) y diferencias pareadas (\(W = 0{,}822\), \(p = 2{,}37 \times 10^{-14}\)). En los tres casos el valor \(p\) queda muy por debajo del nivel de significancia \(\alpha = 0{,}05\), lo que respalda el uso de una prueba no paramétrica. El nivel de significancia se fijó en \(\alpha = 0{,}05\). En lo que sigue se denotan \(T^{+}\) y \(T^{-}\) la suma de los rangos de las diferencias positivas y negativas, respectivamente. La implementación de referencia (\texttt{scipy.stats.wilcoxon} en su variante de dos colas) devuelve como estadístico de contraste el menor de ambos, es decir \(W = \min(T^{+}, T^{-})\); por lo tanto \(W = T^{-}\) cuando el flujo automatizado resulta más rápido que el manual en la mayoría de los pares, que es el caso del corpus analizado. Los \(p\) de tiempos son bilaterales, conservadores frente a la \(H_{1}\) direccional; el unilateral global (\texttt{alternative=greater}) es \(1{,}54 \times 10^{-32}\).

Para cuantificar la magnitud práctica de la diferencia se calculó el coeficiente \(r\) de correlación rank-biserial, indicador de tamaño del efecto para pruebas no paramétricas. El cálculo se define directamente a partir de los dos estadísticos de rango, \(r = (T^{+} - T^{-}) / [n(n+1)/2]\), donde \(T^{+}\) y \(T^{-}\) son las sumas de rangos positivos y negativos y \(n\) el número de pares. La expresión equivale a \(1 - \frac{4W}{n(n+1)}\) cuando \(W = T^{-} = \min(T^{+}, T^{-})\), y su rango teórico es \([-1, +1]\).

Para el análisis de la clasificación se construyó la matriz de confusión global y se calcularon los valores de precisión, sensibilidad (recall) y F1 por cada clase, así como su promedio macro. La referencia de verdad es multietiqueta (107 de los 200 casos, 53,5~\%) y el clasificador emite una etiqueta principal: las métricas primarias son monoetiqueta (exactitud estricta y macro-F1 de umbral) y las secundarias, multietiqueta (subset accuracy, Hamming loss y micro-F1), con su penalización estructural explicitada en la \autoref{sec:resultados}. Como métrica secundaria operativa se calculó además la exactitud laxa (sector predicho principal dentro del conjunto de verdad del caso), distinta de la estricta, que sigue siendo la primaria; ambas se desagregaron por etapa y por canal con intervalo de Wilson, y la confianza se analizó por etapa, porque cada una emite una magnitud distinta. Las proporciones por clase se acotaron con el método de Wilson \parencite{Wilson1927_probable_inference}; el F1 por clase y las métricas de conjunto, con el bootstrap percentil de \textcite{Efron1993_bootstrap} a nivel de caso (2000 remuestreos, semilla fija). La clase de soporte unitario (Bases de Datos, \(n = 1\)) se declara no estimable y se excluye del macro-F1 de umbral, contra el cual se evalúa el criterio de 0,85. Todos los cálculos se realizaron mediante la biblioteca SciPy del ecosistema Python \parencite{Virtanen2020_scipy}.

\subsection{Amenazas a la validez}\label{subsec:amenazas-validez}

Se identificaron y trataron explícitamente cuatro amenazas a la validez del estudio.

La primera, una amenaza a la validez interna por efecto de aprendizaje del operador humano sobre el corpus, se acotó parcialmente mediante la ejecución secuencial pareada descrita en la \autoref{subsec:diseno-experimental}; como se anticipó allí, la rotación de operadores entre las dos mitades no elimina la amenaza (el efecto de operador queda confundido con la mitad del corpus) y se declara como limitación no resuelta del diseño, no como una mitigación efectiva.

La segunda, una amenaza a la validez de constructo por la subjetividad inherente a la categorización de los incidentes, se mitigó mediante el doble etiquetado independiente, el Kappa de Cohen y la resolución de discrepancias por consenso supervisado; su valor (\(\kappa = 0{,}87\)) corresponde a un acuerdo casi perfecto según los rangos de Landis y Koch \parencite{Landis1977_kappa} y respalda la solidez del constructo de medición. Cabe precisar que los analistas etiquetadores son las mismas personas que operaron el brazo manual (\autoref{subsec:poblacion-muestra-corpus}), coincidencia que no invalida el acuerdo inter-evaluador pero acota la independencia del referente y se asume como limitación explícita.

La tercera, una amenaza a la validez externa derivada del idioma y del dominio organizacional, se reconoce como limitación explícita del estudio: los resultados no son extrapolables sin nuevas validaciones a corpus en otros idiomas, en otras organizaciones o en otros sectores de la economía.

La cuarta, una amenaza a la validez de conclusión derivada del tamaño moderado de la muestra (\(n = 200\)), se evaluó mediante un análisis de potencia estadística a priori \parencite{Cohen1988_power_analysis}. Para el contraste principal ---la prueba de Wilcoxon de rangos con signo sobre las mediciones pareadas, bilateral, con \(\alpha = 0{,}05\)--- y una potencia objetivo de \(0{,}80\), el efecto mínimo detectable con \(n = 200\) es un tamaño del efecto estandarizado \(d_z\) (media de las diferencias pareadas sobre su desvío) de aproximadamente \(0{,}20\) en la aproximación normal (\(0{,}21\) con la eficiencia relativa asintótica \(0{,}864\) de la prueba de rangos frente a la \(t\) pareada). El \(d_z\) observado, \(1{,}00\), lo supera más de cuatro veces; por canal es de 1,37 (telefonía), 0,73 (correo) y 1,61 (formulario web), frente a mínimos de 0,31, 0,34 y 0,38 (0,33, 0,37 y 0,41 con el ajuste), de modo que el diseño posee potencia suficiente para el contraste de tiempos. El riesgo se concentra, en cambio, en la precisión de las estimaciones de clasificación (exactitud por clase, macro-F1), que se reportan con intervalos de confianza (Wilson para las proporciones, bootstrap percentil para F1). La amenaza se trata, además, con pruebas no paramétricas robustas frente a tamaños muestrales pequeños y a desviaciones de la normalidad \parencite{HernandezSampieri2018_metodologia}.

\subsection{Gestión del proyecto bajo enfoque Scrumban (método de construcción del artefacto)}\label{subsec:gestion-scrumban}

El presente apartado describe el método de construcción del artefacto, no el método de investigación. El diseño experimental, el muestreo, la construcción del corpus y el análisis estadístico ---esto es, el método del estudio--- quedaron definidos en las Secciones~4.1 a 4.8; la gestión del proyecto bajo Scrumban concierne exclusivamente a cómo se construyó el sistema y se documenta aquí por trazabilidad, sin constituir la metodología de la investigación.

La construcción del sistema se organizó bajo un esquema Scrumban, combinando la planificación por iteraciones propia de Scrum con la gestión continua del flujo de trabajo característica del método Kanban \parencite{Ladas2009_scrumban}. Se ejecutaron seis iteraciones, cuyo detalle y relación con las campañas de medición se presentan en la \autoref{subsec:conexion-scrumban} y en la \autoref{fig:cronograma}.

La distribución de roles asignó la coordinación general, la integración de componentes, el desarrollo del módulo Python y el modelado de la base de datos a Luca Gomez, y la configuración del flujo N8N junto con la integración telefónica mediante Twilio a Gonzalo Sevilla; Carla Bustos participó como revisora de la tesis y colaboradora en los anexos y en la corroboración de las fuentes utilizadas, todo bajo supervisión académica continua del director del trabajo, Prof. Alberto Cortez.

El tablero Kanban se gestiono mediante GitHub Projects, con seguimiento diario del estado de las tareas y reuniones de planificación al inicio de cada iteración. Las definiciones de terminado (\emph{Definition of Done}) incluyeron, para cada tarea, la aprobación de las pruebas unitarias correspondientes, la revisión de código por pares y la actualización de la documentación asociada.

\newpage
